\documentclass[twocolumn]{aastex631}
\begin{document}
\title{The reionization ionizing-photon-budget shortfall for star-forming galaxies is not robust to systematics at z\~{}11 (optimistic systematic corner)}
\author{NebulaMind Lab (autonomous pipeline)}
\affiliation{NebulaMind Lab, \url{https://lab.nebulamind.net}}
\begin{abstract}
Reionization ionizing-photon-budget reconciliation at z\~{}11: star-forming galaxies require f\_esc=0.216 (+0.186/-0.099) to close the budget (Madau-Dickinson SFRD, log xi\_ion=25.5+/-0.15, clumping C=2-5, JWST-SFRD tail), versus indirect-proxy-inferred f\_esc=0.080 (+0.147/-0.051) from LzLCS O32/beta calibrations. Median delta(required-inferred)=+0.118 dex-frac (16-84\%: -0.033 to +0.308); 79\% of the systematic MC shows a shortfall. Conclusion: the budget CLOSES within the systematic, and the sign holds under both O32 and beta calibrations. The result is bounded by the xi\_ion x clumping x proxy-calibration systematic, not by statistics. Generated autonomously from public data (jwst) via the NebulaMind Lab runner: a bounded, reproducible, descriptive study.
\end{abstract}
\section{Introduction}
Introduction:
Recent studies have highlighted a potential crisis in the reionization photon budget, suggesting that star-forming galaxies may not produce enough ionizing photons to account for the observed reionization process [Muñoz2024]. This discrepancy has led to increased scrutiny of the assumptions and parameters used in calculating the ionizing photon budget. To address this issue, we revisit the ionizing photon budget calculation using a literature-anchored approach, building on previous work by Duncan (2015) and Madau (2017). Our aim is to reconcile the reionization photon budget at z\~{}11 and assess whether star-forming galaxies can provide sufficient ionizing photons.
\section{Data and method}
Data and method:
We employ the cosmic SFRD from the Madau \& Dickinson (2014) analytic fitting function, along with published values for xi\_ion and O32/beta f\_esc proxy calibrations [Chisholm+22, Flury+22; Simmonds+24]. Our calculation is based on a systematics reconciliation over these literature values, without relying on new observational or catalog data. We focus specifically on the ionizing-photon-budget method to derive our results.
\section{Result}
Result:
Our analysis reveals that star-forming galaxies require an escape fraction of f\_esc=0.216 (+0.186/-0.099) to close the reionization photon budget at z\~{}11, assuming a Madau-Dickinson SFRD, log xi\_ion=25.5±0.15, and clumping factor C=2-5. In contrast, indirect-proxy-inferred f\_esc values from LzLCS O32/beta calibrations yield f\_esc=0.080 (+0.147/-0.051). The median difference between the required and inferred escape fractions is +0.118 dex-frac (16-84\%: -0.033 to +0.308), with 79\% of systematic Monte Carlo simulations showing a shortfall in ionizing photons.
\begin{figure}\centering\includegraphics[width=\columnwidth]{result.png}\caption{The reionization ionizing-photon-budget shortfall for star-forming galaxies is not robust to systematics at z\~{}11 (optimistic systematic corner)}\end{figure}
\section{Caveats}
Caveats:
Our study relies on an automated, single-selection, uncalibrated measurement approach, which has inherent limitations. The accuracy of our results depends heavily on the assumptions and calibrations used in previous studies, such as the Madau \& Dickinson (2014) SFRD and O32/beta f\_esc proxy calibrations. Additionally, our analysis does not account for potential systematic errors or uncertainties in these underlying parameters. Furthermore, the use of a fixed clumping factor may oversimplify the complex reionization process, which could affect the accuracy of our results. Therefore, while our study provides valuable insights into the reionization photon budget crisis, it should be interpreted with caution and considered alongside other complementary research approaches.
\section*{References}
\footnotesize
[Muoz2024] Muñoz et al. (2024). Reionization after JWST: a photon budget crisis?. bibcode 2024MNRAS.535L..37M \par [Park2022] Park et al. (2022). Calibrating excursion set reionization models to approximately conserve ionizing photons. bibcode 2022MNRAS.517..192P \par [Duncan2015] Duncan et al. (2015). Powering reionization: assessing the galaxy ionizing photon budget at z < 10. bibcode 2015MNRAS.451.2030D \par [Davies2021] Davies et al. (2021). The Predicament of Absorption-dominated Reionization: Increased Demands on Ionizing Sources. bibcode 2021ApJ...918L..35D \par [Madau2017] Madau (2017). Cosmic Reionization after Planck and before JWST: An Analytic Approach. bibcode 2017ApJ...851...50M

\end{document}
